% =====================================================================
%  short6.tex -- 6-page condensed version
%  "Layered, Evidence-Aware Bluetooth Fingerprinting"
%
%  SCOPE RULE for this file: it contains ONLY claims, results and figures
%  that appear in the project report (Irem's_Report_EN). Nothing from the
%  longer latest.tex that the report does not state has been carried over.
%  In particular this draft deliberately omits: the repeat-scan stability
%  study, the three implementation defects found by the new pair, the
%  chronology and vendor-text checks, confidence intervals, and the
%  GDPR/ethics section. Those live in latest.tex.
%  EXCEPTION 2, requested separately: the worked examples of Section III-F
%  (Table~\ref{tab:examples}). They are replayed from the archived scans with
%  research/replay_evidence_chain.py against the current scorer, so they are
%  project output rather than report content, and one of them states a
%  limitation the report does not contain (see the text).
%  EXCEPTION, requested separately: the capability comparison table
%  (Table~\ref{tab:cap}) and its section. It is condensed from tab:signals
%  and tab:cap in latest.tex, restricted to the method families the report
%  itself names, and kept qualitative so that no measured figure absent from
%  the report is reintroduced.
%
%  STATUS: draft. Resolve author block before submission.
% =====================================================================
\documentclass[conference]{IEEEtran}
\IEEEoverridecommandlockouts

\usepackage{cite}
\usepackage{array}
\usepackage{booktabs}
\usepackage{url}
\usepackage[T1]{fontenc}

\begin{document}

\title{Layered, Evidence-Aware Bluetooth Fingerprinting:\\
Low-Cost Signals and Conflict-Aware Inference}

% TODO before submission: if a second author is added, restore the \and block
% that latest.tex carries and confirm the ordering with the supervisor.
\author{\IEEEauthorblockN{Irem Cesur}
\IEEEauthorblockA{CYSECDIGITAL Research Group\\
Technische Hochschule Mittelhessen (THM)\\
Friedberg, Germany}}

\maketitle

\begin{abstract}
We infer the type and operating system of a nearby Bluetooth device without
changing that device's state. Knowing the operating system supports security
auditing, because many documented weaknesses are specific to particular
versions, so an OS label narrows which of them could apply. Existing methods
each reach part of this: they flood the target, read a single protocol field,
or only listen to the radio, and each falls short in some respect. We aim
instead for three properties together --- low interaction cost, the ability to
distinguish many operating systems, and reporting conflicts between signals
rather than hiding them. We extend an existing baseline scanner with
additional signal layers obtained at low interaction cost, add one cache-read
signal usable on re-encounter, and attach an \emph{evidence chain} that leaves
the prediction unchanged but labels each signal relative to it and lowers the
reported confidence on a genuine conflict. We then measure how far that
mechanism works on real data, which produced our main finding. Both signal
pairs the chain originally compared depended on GATT~Appearance, present in 2
of 149 records and 0 of 77 scans, so the chain never fired; a pair requiring
no Appearance raised coverage from 0\% to 48\% of scans, but to 62\% on our
own devices against 10\% on third-party ones. What limits this class of
inference is therefore not the quality of a signal but the simultaneous
availability of the signals a check requires. We claim no calibrated accuracy
result, and present this as a measurement that openly reports a constraint we
observed.
\end{abstract}

\begin{IEEEkeywords}
Bluetooth fingerprinting, BLE privacy, operating-system identification,
evidence-based inference, uncertainty, security audit.
\end{IEEEkeywords}

% =====================================================================
\section{Introduction}
% =====================================================================
Bluetooth is present in phones, wearables, peripherals, and IoT devices
almost everywhere. The focus of this work is to infer the type and operating
system of a nearby Bluetooth device without changing the device's state. We
build on a baseline active scanner developed earlier within our group.

The value of doing so is auditing. If we know a device's operating system, we
can audit its security, and since various vulnerabilities are specific to
certain versions, an OS indicator considerably narrows down which
vulnerabilities of that device could be
exploited~\cite{barua2022blethreats}. Bluetooth was never designed to expose
such an identifier, and privacy mechanisms such as MAC-address randomization
work against recovering one~\cite{martin2017macrandom}.

Existing detection methods almost do this job, but each of them falls short in
some respect. Flooding-based probing recovers an OS version by observing the
target under stress~\cite{kavisankar2024osid}, at a high interaction cost and
with no representation of uncertainty. Methods that read a single protocol
field or a set of service metadata are cheap~\cite{celosia2019gatt} but stop
short of an operating system. Methods that only listen to the radio are
passive and hard to spoof~\cite{santanacruz2024jsdivergence} but need special
hardware and return an identity rather than an OS.

This work aims for three desirable properties together: low interaction cost,
the ability to distinguish many operating systems, and reporting conflicts
between signals rather than hiding them. In other words, rather than producing
a new signal, the aim is both to enrich what is already available and to make
uncertainties visible.

Our four contributions are:

\begin{itemize}
  \item \textbf{Added signal layers at low interaction cost}
  (Section~\ref{sec:layers}): the LMP feature bitmap, chipset and
  manufacturer resolution, GATT~Appearance, GATT connection parameters, and
  the full SDP attribute tree.
  \item \textbf{A cache-read signal usable on re-encounter}
  (Section~\ref{sec:layers}): Modalias, read from the host cache rather than
  from the device.
  \item \textbf{An evidence chain for conflict-aware inference}
  (Section~\ref{sec:evidence}): each signal is labelled relative to the
  inference, and genuine disagreement is reported rather than averaged away.
  \item \textbf{A measurement of how, and to what extent, the mechanism works
  on real data} (Section~\ref{sec:eval}).
\end{itemize}

\noindent\textbf{Constraints and non-goals.} The most fundamental constraint
is that every signal this method uses --- Class-of-Device, device name,
GATT~Appearance, Modalias, SDP --- is declared by the target device and is not
cryptographically bound. The device may declare them falsely and mislead us,
and we know of no existing method to detect this. In addition, MAC
randomization causes the same device to be seen as multiple addresses; our
sample is small, so the rates we measure carry wide confidence intervals; and
the confidence score we produce is not calibrated. Accordingly, this work
infers a device type and an operating-system family, but does not claim to
extract an exact OS version number, does not break MAC randomization, and does
not perform an accuracy comparison with other methods on shared data. Such a
comparison would require a ground-truth dataset, a layer-by-layer ablation,
and a calibrated confidence function. What we actually measure here,
separately from all of these, is the availability of the signals.

% =====================================================================
\section{Background}\label{sec:related}
% =====================================================================
It is worth being precise about what the existing families of work do, since
the contribution here is not a better signal but a fuller use of the ones
already within reach.

\textbf{Flooding the target: Kavisankar \emph{et al.}}~\cite{kavisankar2024osid}
The closest work in intent identifies the operating system of a
Bluetooth device by sending traffic and observing how the device responds.
This reaches a fine-grained answer, and it pays for it twice: the target's
state is disturbed, which is exactly what we set out to avoid, and the result
arrives as a single answer with no account of how well supported it was.

\textbf{Reading a single protocol field: Celosia and Cunche}~\cite{celosia2019gatt}
Fingerprinting from the Generic Attribute Profile a BLE device exposes shows
that declared metadata alone is often enough to name a device. This is cheap and
does not disturb the device, and we borrow from it directly. Its limit is
scope: a device model is not an operating-system family, and a method resting
on one field inherits whatever that field happens to contain.

\textbf{Only listening to the radio: Santana-Cruz \emph{et al.}}~\cite{santanacruz2024jsdivergence}
Identifying a device from imperfections in its radio emissions is the
hardest of the three to spoof, because the device is never asked to declare
anything. It is also the least deployable, requiring dedicated hardware, and
it recognises \emph{that} device again rather than reporting what it runs.

\textbf{Our position.} We sit with the declarative methods and accept their
central weakness, stated above: every signal is declared by the target. What
we add is breadth, several declared signals rather than one, and an explicit
account of what happens when those signals do not agree with each other.

% =====================================================================
\section{Approach}\label{sec:approach}
% =====================================================================
\subsection{Layered structure}\label{sec:layers}
The system scans a single target address in layers. Each layer is defined by
what it costs to obtain rather than by what it reveals.

\textbf{Layer~0: the baseline active scan.} The inherited scanner, which reads
the Class-of-Device, the SDP service list, the GATT service identifiers, the
vendor prefix of the address (OUI), and the LMP protocol version.

\textbf{Layer~1: parsing what is already there, and one cache read.} The
baseline already issues the query whose response carries the LMP feature
bitmap; we parse that bitmap rather than discarding it, so the signal costs
parsing and no new connection. Layer~1 also reads \emph{Modalias} from the
BlueZ cache, and resolves chipset and manufacturer from it where present.
Reading the cache requires no live connection at scan time. It does require
that the host has previously interacted with the device, so Modalias is a
signal for a re-encounter or a re-scan rather than a generally free one. We
state this as a scope limit here, and Section~\ref{sec:eval} confirms it
empirically rather than assuming it.

\textbf{Layer~2: BLE category and the full SDP tree.} We read
GATT~Appearance, which matters especially for BLE-only devices that expose no
Class-of-Device, together with the GATT connection parameters; and where the
baseline reads only a summarised service list, we read the full SDP attribute
tree.

One point here matters more than its technical detail suggests: \emph{these
signals do not share a single connection}. SDP goes over the BR/EDR connection
the baseline already opens, but GATT is read over a separate LE transport. A
device that answers on one may be silent on the other, and
Table~\ref{tab:transport} gives the access path per signal for that reason. It
is the table that explains most of the availability results later.

\begin{table}[t]
\centering
\footnotesize
\renewcommand{\arraystretch}{1.15}
\caption{Access path per signal. The BR/EDR signals reuse the connection the
baseline already opens; GATT is read over a separate LE transport; the cache
read needs no live connection at scan time.}
\label{tab:transport}
\begin{tabular}{@{}ll@{}}
\toprule
\textbf{Signal} & \textbf{Access path} \\
\midrule
LMP version and feature bitmap & BR/EDR (reuses baseline) \\
SDP service list and attribute tree & BR/EDR (reuses baseline) \\
Class-of-Device & inquiry / host cache \\
GATT Appearance, conn.\ params, UUIDs & LE (separate transport) \\
Modalias & host cache (no live conn.) \\
RSSI / RTT (Layer~3) & active observation window \\
\bottomrule
\end{tabular}
\end{table}

\textbf{Layer~3: time-windowed behavioural observation.} A short windowed
observation such as RSSI and RTT. It is measured and reported, but it does not
feed back into the current inference. We state this explicitly because it
bounds every claim that follows: no inference in this work relies on this
layer.

\subsection{Evidence chain}\label{sec:evidence}
The methodological core of this work is not a signal but a discipline about
how signals are combined.

The evidence chain is an audit layer produced separately, after the inference
has already been determined. It does not change the inference. What it does is
label each signal as \emph{primary}, \emph{supports}, \emph{conflicts} or
\emph{unused} relative to that inference, and lower the reported confidence in
the case of a genuine conflict. This is the point of departure from a single
fused score: if two signals point in opposite directions, a score records only
their sum, and the operator never learns that the device contradicted itself.

\textbf{Avoiding circularity.} The \emph{primary} label carries the idea that
makes the rest work. If an inference was produced by a signal, that signal is
labelled \emph{primary} and excluded from the count, because comparing it with
the result it derived itself would be self-validation. Only signals that
played no part in producing the inference are eligible to support or
contradict it. The same argument recurs throughout this design, and it is the
reason several things one might expect to find here are deliberately absent.

\textbf{Two stages.} The process takes place in two stages. First, each signal
is evaluated against the inference. Second, pairs of independently produced
signals are evaluated against each other rather than against the inference. A
pair is evaluated only when both of its halves resolved: a device on which
only one half is present yields no record, since an absent check is not a
passed check. That property is unremarkable in itself, and it turns out to
govern the central result of Section~\ref{sec:eval}.

\subsection{Re-encounter comparison}\label{sec:reencounter}
When a device is scanned twice, the question that comes to mind is whether
this scan is consistent with the previous one. The obvious way of answering it
--- comparing the two inferences --- is closed to us by the circularity
argument above. Agreement between two inferences measures not accuracy but the
system's internal consistency, because both inferences come from the same
scorer.

Therefore, in this part we compare only the invariant fields: Modalias,
Class-of-Device, and the OUI-resolved manufacturer. These fields cannot differ
within a short interval; if they differ, there is an error in at least one
place, and that conclusion is available without any ground truth. A
disagreement is recorded as \emph{conflicts}. An agreement is counted as
\emph{unused}, because a reader agreeing with itself is not independent
corroboration. A device scanned for the first time does not enter this
comparison at all.

\subsection{Scoring model}\label{sec:scoring}
The device type is determined by a branching system with a fixed priority
order, in which the first match wins: Apple, Android, embedded IoT,
PC/laptop, wearable, audio. Only the Android branch has a numerical model; the
remaining branches are rule-based. We describe the model at this level
deliberately, since the weights it uses are not calibrated
(Section~\ref{sec:approach}) and presenting them as tuned parameters would
overstate what they are.

\subsection{Worked examples}\label{sec:examples}
The mechanisms above are easier to judge against concrete records than in the
abstract. Table~\ref{tab:examples} lists every distinct device in the full-scan
archive --- 26 devices over 87 scans --- replayed against the current scorer,
and groups them by what the evidence chain was able to do with each.

The grouping is itself the result. The pair of Section~\ref{sec:evidence} could
be evaluated on 9 of the 26 devices; on the remaining 17 at least one half did
not resolve, so no record was produced and the chain was empty. Rows~1 and~11
show why that matters: they are the same physical earbuds, with the same cached
Modalias and the same Class-of-Device, distinguished only by whether the
advertised name resolved. With the name, the cascade returns an audio
peripheral running an embedded firmware; without it, the same device is
reported as an Apple phone or computer running iOS~17+. Both answers carry
confidence \emph{High}, and in the second the chain has nothing to say, because
an absent check is not a passed check.

Rows~8 and~9 state a limitation of the two-stage design that we had not
identified before assembling this table. On the smartwatch and on the Samsung
earbuds, the two category signals agree with each other --- wearable and
wearable, audio and audio --- so each pair is recorded as \emph{supports}. The
prediction, however, is \emph{Smartphone} in both cases, contradicting both
halves of the pair it was just told agree. Nothing in the chain catches this:
the second stage compares the halves \emph{to each other} and never to the
prediction, while the first stage does not cover these signals at all, since
they fed the cascade and are therefore \emph{primary} and excluded by the
circularity rule of Section~\ref{sec:evidence}. The rule exists so that a
signal cannot corroborate a prediction it produced; in excluding those signals
it also gives up the ability to notice that the prediction disagrees with all
of them at once. A third check --- comparing an agreeing pair against the
prediction and flagging the case where it matches neither half --- would catch
both rows, and we state it as future work rather than as something this system
does.

Two further patterns in the table bear on the availability argument of
Section~\ref{sec:eval}, and neither is flattering. First, rows~12 and~16: the
OUI resolves the vendor as Apple, the cached Modalias resolves the vendor as
Apple, and the prediction is Android~15 at confidence \emph{High}. Here the
contradicting evidence \emph{was} available and no stage of the chain looks at
it, because no check compares a resolved vendor against a predicted OS
family. That is a missing check rather than a missing signal, and it is a
different failure from the co-availability one this paper argues for.

Second, the largest group in the table is the 13 devices that advertised a name
which our derivation mapped to no category. Several of those names ---
\emph{realme 11 Pro+ 5G}, \emph{OPPO A5 2020}, and two Samsung flagships named
after their model --- are perfectly serviceable product strings. The pair did
not fire on them because our keyword table did not cover them, not because the
device withheld anything. Co-availability, as we measure it, is therefore
partly a property of the environment and partly an artefact of our own
derivation, and the 48\% coverage figure of Section~\ref{sec:eval} should be
read with that split in mind.

\begin{table*}[!t]
\centering
\scriptsize
\setlength{\tabcolsep}{3.4pt}
\renewcommand{\arraystretch}{1.18}
\caption{Every distinct device in the full-scan archive (26 devices, 87 scans),
replayed against the current scorer and grouped by what the evidence chain was
able to do. \textbf{CoD} is the Class-of-Device major class; \textbf{name cat.}
is the major class implied by the advertised name; the pair of
Section~\ref{sec:evidence} is evaluated only when both resolve. \textbf{Mod.}
is the Modalias vendor and product prefix, \textbf{LMP} the link-layer version.
Names chosen by their owners are redacted as $\langle$name$\rangle$; the rest
are product strings the devices advertise. Predictions in bold are those where
the vendor is resolved as Apple and the predicted OS is Android. Every
prediction here is reported at confidence \emph{High} except the two
\emph{Unknown Peripheral} rows, which are \emph{Low}.}
\label{tab:examples}
\begin{tabular}{@{}c
>{\raggedright\arraybackslash}p{2.62cm} c c c c c c
>{\raggedright\arraybackslash}p{4.15cm} c@{}}
\toprule
& \textbf{Advertised name} & \textbf{Scans} & \textbf{CoD} & \textbf{name cat.}
& \textbf{Mod.} & \textbf{LMP} & \textbf{Vendor} & \textbf{Prediction}
& \textbf{Chain} \\
\midrule
\multicolumn{9}{@{}l}{\emph{Pair evaluated; prediction consistent with it}} \\
\midrule
1 & $\langle$name$\rangle$ (AirPods) & 13 & \texttt{0x04} & audio & \texttt{v004Cp201} & 14 & Apple & Audio Peripheral (Apple), Embedded Firmware (RTOS) \newline \emph{Audio Sink Device (LMP 14)} & supports \\
2 & $\langle$name$\rangle$ iPhone & 11 & \texttt{0x02} & phone & \texttt{v004Cp741} & 12 & Apple & Smartphone/Computer (Apple), iOS / macOS \newline \emph{iOS 16.x / macOS 13.x} & supports \\
3 & Redmi A5 & 2 & \texttt{0x02} & phone & \texttt{v00E0p000} & 11 & Xiaomi & Smartphone, Android \newline \emph{Android 15 (Vanilla Ice Cream)} & supports \\
4 & $\langle$name$\rangle$.iPhone & 1 & \texttt{0x02} & phone & \texttt{v004Cp741} & 11 & Apple & Smartphone (iPhone), iOS \newline \emph{iOS 15{+}} & supports \\
5 & iPad von $\langle$name$\rangle$ & 1 & \texttt{0x01} & computer & \texttt{v004Cp770} & 12 & Apple & Smartphone/Computer (Apple), iOS / macOS \newline \emph{iOS 16.x / macOS 13.x} & supports \\
6 & iPad (2) & 1 & \texttt{0x01} & computer & \texttt{v004Cp790} & 12 & Apple & Smartphone/Computer (Apple), iOS / macOS \newline \emph{iOS 16.x / macOS 13.x} & supports \\
7 & $\langle$name$\rangle$ \.{I}pad & 1 & \texttt{0x01} & computer & \texttt{v004Cp761} & 12 & Apple & Computer (Mac), macOS \newline \emph{macOS 13{+}} & supports \\
\midrule
\multicolumn{9}{@{}l}{\emph{Pair evaluated; \textbf{prediction contradicts both halves of the pair}}} \\
\midrule
8 & HUAWEI WATCH FIT 3-8CF & 8 & \texttt{0x07} & wearable & -- & 11 & Huawei & Smartphone, Android \newline \emph{Android 13 (Tiramisu)} & supports \\
9 & Galaxy Buds2 (C7D2) & 5 & \texttt{0x04} & audio & \texttt{v0075pA01} & 11 & Samsung & Smartphone, Android \newline \emph{Android 15 (Vanilla Ice Cream)} & supports \\
\midrule
\multicolumn{9}{@{}l}{\emph{No pair: Class-of-Device present, advertised name resolved to no category}} \\
\midrule
10 & realme 11 Pro{+} 5G & 12 & \texttt{0x02} & -- & \texttt{v00E0p000} & 12 & Realme & Smartphone, Android \newline \emph{Android 15 (Vanilla Ice Cream)} & \emph{empty} \\
11 & \emph{(none)} & 5 & \texttt{0x04} & -- & \texttt{v004Cp201} & 14 & Apple & Smartphone/Computer (Apple), iOS / macOS \newline \emph{iOS 17{+} / macOS 14{+}} & \emph{empty} \\
12 & $\langle$name$\rangle$ & 4 & \texttt{0x02} & -- & \texttt{v004Cp750} & 12 & Apple & \textbf{Smartphone, Android} \newline \emph{Android 15 (Vanilla Ice Cream)} & \emph{empty} \\
13 & OPPO A5 2020 & 2 & \texttt{0x02} & -- & \texttt{v001Dp120} & 9 & OPPO & Smartphone, Android \newline \emph{Android 14 (Upside Down Cake)} & \emph{empty} \\
14 & \emph{(none)} & 2 & \texttt{0x04} & -- & \texttt{v0075pA01} & -- & Samsung & Unknown Peripheral, Unknown & \emph{empty} \\
15 & \emph{(none)} & 1 & \texttt{0x02} & -- & \texttt{v00E0p000} & -- & Realme & Unknown Peripheral, Unknown & \emph{empty} \\
16 & $\langle$name$\rangle$ & 1 & \texttt{0x02} & -- & \texttt{v004Cp760} & 12 & Apple & \textbf{Smartphone, Android} \newline \emph{Android 15 (Vanilla Ice Cream)} & \emph{empty} \\
17 & S22 Ultra von $\langle$name$\rangle$ & 1 & \texttt{0x02} & -- & \texttt{v0075p010} & 11 & Samsung & Smartphone, Android \newline \emph{Android 15 (Vanilla Ice Cream)} & \emph{empty} \\
18 & S25{+} von $\langle$name$\rangle$ & 1 & \texttt{0x02} & -- & \texttt{v0075p010} & 13 & Samsung & Smartphone, Android \newline \emph{Android 15 (Vanilla Ice Cream)} & \emph{empty} \\
19 & \emph{(none)} & 1 & \texttt{0x01} & -- & \texttt{v0006p000} & 8 & Intel & PC / Laptop, Windows \newline \emph{Windows 8.1 / 10} & \emph{empty} \\
20 & $\langle$name$\rangle$ & 1 & \texttt{0x01} & -- & \texttt{v0006p000} & 13 & Liteon & PC / Laptop, Windows \newline \emph{Windows 11} & \emph{empty} \\
21 & ``a (2)'' & 1 & \texttt{0x02} & -- & -- & 14 & Apple & Smartphone (iPhone), iOS \newline \emph{iOS 17{+}} & \emph{empty} \\
22 & $\langle$name$\rangle$ & 1 & \texttt{0x02} & -- & \texttt{v004Cp761} & 12 & Apple & Smartphone (iPhone), iOS \newline \emph{iOS 16{+}} & \emph{empty} \\
\midrule
\multicolumn{9}{@{}l}{\emph{No pair: name category present, no Class-of-Device}} \\
\midrule
23 & $\langle$name$\rangle$-LAPTOP & 1 & -- & computer & -- & 9 & Fugui & PC / Laptop, Windows \newline \emph{Windows 8.1 / 10} & \emph{empty} \\
\midrule
\multicolumn{9}{@{}l}{\emph{No pair: neither half resolved}} \\
\midrule
24 & \emph{(none)} & 8 & -- & -- & -- & 13 & Liteon & PC / Laptop, Windows \newline \emph{Windows 11} & \emph{empty} \\
25 & $\langle$name$\rangle$ & 1 & -- & -- & -- & 13 & Liteon & PC / Laptop, Windows \newline \emph{Windows 11} & \emph{empty} \\
26 & $\langle$name$\rangle$'s S22 & 1 & -- & -- & -- & 12 & Samsung & Smartphone, Android \newline \emph{Android 15 (Vanilla Ice Cream)} & \emph{empty} \\
\bottomrule
\end{tabular}
\end{table*}

\subsection{Differences from the baseline}
In the baseline work, every signal contributed to a single weighted score and
a single decision was made without uncertainty, and Modalias was not read. We
added two things: the new signal layers of Section~\ref{sec:layers} and the
evidence chain of Section~\ref{sec:evidence}. The result is not a new
detector. With almost zero new connections, it uses the resources already at
hand and turns the system into one in which disagreement between signals is
reported.

% =====================================================================
\section{Comparison with Existing Methods}\label{sec:comparison}
% =====================================================================
The methods of Section~\ref{sec:related} are evaluated on different device
sets with different mechanisms, so their reported accuracies are not
comparable and we do not tabulate them as if they were; as stated in the
introduction, an accuracy comparison on shared data is explicitly not
something this work performs. What can be compared is what each method
\emph{requires}, what it \emph{emits}, and how it behaves when its evidence is
weak. Table~\ref{tab:cap} sets that out along the three properties this work
set out to obtain together.

Read across the rows, the table states the claim of the introduction in a
compact form. Each existing family obtains two of the three properties at
most. The flooding method reaches an OS version but disturbs the target and
reports one verdict whatever the evidence behind it; the declarative metadata
methods are cheap and quiet but stop at a device model; the radio methods ask
the device to declare nothing at all, which makes them the hardest to
mislead, but they need dedicated hardware and return an identity rather than
an operating system. The baseline we extend reaches an OS but folds every
signal into one number, so a disagreement between two of its inputs leaves no
trace in its output.

The last row is not a claim of superiority, and one column is deliberately
unflattering. This work obtains the three properties together, but the
conflict reporting it adds is only as available as the signals a check needs
--- which is what Section~\ref{sec:eval} measures, and where it arrives at
48\% of scans overall and 10\% on third-party devices. A mark in the final
column therefore describes a mechanism that exists, not one that fires on
every scan.

\begin{table*}[!t]
\centering
\footnotesize
\setlength{\tabcolsep}{5pt}
\renewcommand{\arraystretch}{1.3}
\caption{Which properties each method has. The three rightmost
columns are the properties this work set out to obtain together (Section~I).
Cells are qualitative: the methods report results on disjoint device sets, so
no row is comparable to another as a score, and none is given here. ``Interaction
cost'' means radio activity beyond passive discovery. This work is the last
row.}
\label{tab:cap}
\begin{tabular}{@{}>{\raggedright\arraybackslash}p{3.3cm}
>{\raggedright\arraybackslash}p{2.9cm}
>{\raggedright\arraybackslash}p{3.5cm}
>{\raggedright\arraybackslash}p{2.6cm}
>{\raggedright\arraybackslash}p{3.4cm}@{}}
\toprule
\textbf{Method} & \textbf{Signals it reads} & \textbf{Low interaction cost}
& \textbf{Distinguishes many OSes} & \textbf{Reports conflicts between signals} \\
\midrule
Kavisankar \emph{et al.}~\cite{kavisankar2024osid}\newline \emph{(flooding the target)} &
Responses under induced stress; timing &
\textbf{No}: floods the target and changes its state &
\textbf{Yes}: OS and version &
\textbf{No}: one verdict, always emitted \\
\addlinespace
Celosia and Cunche~\cite{celosia2019gatt}\newline \emph{(single protocol field)} &
GATT attribute profile &
\textbf{Yes}: one read, device undisturbed &
\textbf{No}: device model, not an OS &
\textbf{No}: a uniqueness score, not a disagreement \\
\addlinespace
Santana-Cruz \emph{et al.}~\cite{santanacruz2024jsdivergence}\newline \emph{(radio layer only)} &
Radio-layer emissions &
\textbf{Yes} in radio terms, but requires dedicated hardware &
\textbf{No}: a device identity, not an OS &
\textbf{No}: a confidence value only \\
\addlinespace
% TODO: the report attributes the baseline to a former group member (Aniket).
% If it should be named or cited here, supply the surname and a citable
% artifact (thesis, internal report, repository); do not leave it anonymous
% by default.
Baseline scanner, this group\newline \emph{(extended here)} &
CoD, SDP list, GATT UUIDs, OUI, LMP version &
Medium: one connect-and-query pass &
\textbf{Yes}: OS and version &
\textbf{No}: all signals fused into one weighted score \\
\addlinespace
\textbf{This work} &
The baseline's signals, plus the LMP feature bitmap, Modalias, GATT
Appearance and connection parameters, and the full SDP tree &
\textbf{Yes}: Layers~1--2 reuse the baseline's connection or read the host
cache; almost zero new connections &
\textbf{Yes}: OS family and version, same branches as the baseline &
\textbf{Yes}: explicit conflict flag and lowered confidence --- but gated on
signal co-availability (Section~\ref{sec:eval}) \\
\bottomrule
\end{tabular}
\end{table*}

% =====================================================================
\section{Preliminary Signal-Availability Measurement}\label{sec:eval}
% =====================================================================
The question this section answers is how often the Layer~1--2 signals we added
are actually present in the surroundings. This is not an accuracy study, and
we treat it as the prior question rather than a lesser one: a rarely present
signal, no matter how well it classifies, cannot help us as long as it is not
present. In this sense availability took precedence over accuracy.

\textbf{Method.} A standalone audit passively discovers Bluetooth addresses
and then runs only the Layer~1--2 queries against each, recording per device
whether it was \emph{reached} at all and, if so, which signals it exposed. We
separate availability from the signal being rich: a device that opens a
connection but exposes no tracked signal is counted as reached, not as a
connection failure, since conflating the two would make the fraction of
reached devices exposing a given signal partly circular. Two availability
rounds were run: a pilot in a single office and a wider sweep across several
locations. Separately, a set of full scans was archived, and those archives
are what the evidence-chain measurement below replays.

\textbf{Reachability is the first constraint.} The pipeline could not reach
the vast majority of addresses at all: 16 of 18 in the pilot, and 46 of 52 in
the sweep (Table~\ref{tab:avail}). Whatever the added signals are worth, they
are worth nothing on an address the pipeline never opens a channel to, and
this dominates every other availability figure we report.

\begin{table}[t]
\centering
\footnotesize
\setlength{\tabcolsep}{4pt}
\renewcommand{\arraystretch}{1.2}
\caption{Signal availability in the wild. Modalias is a host-cache read rather
than a live signal, and is reported against all addresses in the round. Counts
are given raw: at these sample sizes a percentage would read as an estimate it
cannot support.}
\label{tab:avail}
\begin{tabular}{@{}>{\raggedright\arraybackslash}p{4.7cm} r r@{}}
\toprule
& \textbf{Pilot ($n{=}18$)} & \textbf{Sweep ($n{=}52$)} \\
\midrule
Never reached by pipeline & 16 & 46 \\
Reached & 2 & 6 \\
\midrule
Modalias present (of all addresses) & 8/18 & 0/52 \\
\bottomrule
\end{tabular}
\end{table}

\textbf{Modalias is a re-scan signal, and the measurement shows it.} Modalias
came out at 44\% on previously paired devices in the office, while it was not
found at all in scans at different locations. The pilot ran where the host had
already dealt with those devices; the sweep met unfamiliar ones and found
nothing cached. This is the scope claim of Section~\ref{sec:layers} turned
into evidence rather than left as an assurance: Modalias is only useful on
re-scan.

\textbf{The evidence chain did not work on real data, and the reason is the
finding.} Having built the chain, we examined how it works on the archived
scans, and found that it did not fire. The cause was not in its logic. Both of
the first two pairs we had chosen depended on GATT~Appearance, and this signal
was present in only 2 of 149 records and in 0 of 77 scans. We had gated the
entire mechanism on the scarcest signal in the study.

The error is instructive because of how it was made. We had chosen the pairs
according to how \emph{independent} they were, and had not paid attention to
how often they are present \emph{together}. Choosing a new pair that does not
require Appearance --- Class-of-Device against the advertised name --- raised
coverage from 0\% to 48\% of scans.

\begin{table}[t]
\centering
\footnotesize
\setlength{\tabcolsep}{5pt}
\renewcommand{\arraystretch}{1.15}
\caption{Coverage of the Class-of-Device versus advertised-name pair, split by
whose devices were scanned. Coverage means both halves resolved, so the pair
could be evaluated at all.}
\label{tab:pairsplit}
\begin{tabular}{@{}>{\raggedright\arraybackslash}p{4.2cm} r@{}}
\toprule
\textbf{Population} & \textbf{Coverage} \\
\midrule
Own devices & 62\% \\
Third-party devices & 10\% \\
\midrule
All scans & 48\% \\
\bottomrule
\end{tabular}
\end{table}

\textbf{Coverage does not transfer to unfamiliar devices.}
Table~\ref{tab:pairsplit} carries the result we consider most important, and
it is the least flattering one here. The method worked at 62\% on our own
devices and at only 10\% on third-party devices. The own-device figure is
therefore an upper bound that does not transfer to the setting the method
exists for, and reporting it alone would have overstated the method's reach by
roughly a factor of six.

\textbf{What this means for the design.} What limits this kind of inference is
not the quality of the signal but the simultaneous availability
--- the co-availability --- of the required signals. A pair is useful only if
both halves are also commonly present, and we can draw a further conclusion
from this: independence and availability work against each other. The signals
that are most clearly independent of the primary decision path are precisely
the ones that path did not need, and therefore the ones a device is least
likely to expose. At our sample sizes we state this as the constraint that
governs how such a mechanism should be built, and as a hypothesis the pair
change above is designed to test, rather than as a settled result.

\textbf{Limitations of this measurement.} It measures availability, not
identification accuracy. The sample is small, so every proportion carries a
wide confidence interval, and the rounds count \emph{addresses} rather than
distinct devices, which MAC randomization means is not the same thing. The
evidence chain, though it now fires, has not yet produced a genuine
device-originated conflict on real data, and Section~\ref{sec:examples} shows
a class of disagreement --- the prediction contradicting an agreeing pair ---
that its current two stages cannot express at all.

% =====================================================================
\section{Conclusion}\label{sec:conclusion}
% =====================================================================
In this work we added low-cost new signal layers, and an evidence chain that
makes conflicts visible, to an existing baseline scanner. The goal was not so
much to produce a new signal as to enrich the resources at hand and report
uncertainty honestly.

When we measured the mechanism on real data, we learned the actual lesson.
What limits this kind of inference is not the quality of the signal but the
co-availability of the required signals: the method worked at 62\% on our own
devices and only 10\% on third-party ones, and before that it had not worked
at all, because both of its original pairs rested on a signal present in 2 of
149 records. We therefore present this work not as an accuracy success but as
a measurement that openly reports a constraint we observed.

The next steps follow directly from the non-goals stated in the introduction.
An accuracy claim needs a ground-truth dataset, a layer-by-layer ablation, and
a calibrated confidence function; each of those is future work, and none of
them is assumed here.

% =====================================================================
\section*{Acknowledgment}
% =====================================================================
This work was carried out at the CYSECDIGITAL research group, THM.
% TODO before submission: add advisor, collaborators, and any funding.

% =====================================================================
%  REFERENCES --- re-verify every entry against the primary source.
%  Only the works the report itself refers to are cited: the three method
%  families it names, the version-specific-vulnerability claim, and MAC
%  randomization.
% =====================================================================
\begin{thebibliography}{00}

\bibitem{kavisankar2024osid}
L.~Kavisankar, A.~Vemuri, S.~Venkatesan, and R.~Khondoker,
``A Novel Strategy for the Identification of the Operating System of
Bluetooth-Enabled Devices for Security Audit,''
in \emph{Proc.\ 11th Int.\ Conf.\ Internet of Things: Systems,
Management and Security (IOTSMS)}, IEEE, 2024.

\bibitem{celosia2019gatt}
G.~Celosia and M.~Cunche,
``Fingerprinting Bluetooth-Low-Energy Devices Based on the Generic
Attribute Profile,''
in \emph{Proc.\ 2nd Int.\ ACM Workshop on Security and Privacy for the
Internet-of-Things (IoT S\&P)}, 2019.

\bibitem{santanacruz2024jsdivergence}
R.~F.~Santana-Cruz, M.~Moreno-Guzman, C.~E.~Rojas-L\'opez,
R.~V\'azquez-Mor\'an, and R.~V\'azquez-Medina,
``Bluetooth Device Identification Using RF Fingerprinting and
Jensen--Shannon Divergence,''
\emph{Sensors}, vol.~24, no.~5, art.~1482, 2024.

\bibitem{martin2017macrandom}
J.~Martin \emph{et al.},
``A Study of MAC Address Randomization in Mobile Devices and When It
Fails,''
\emph{Proc.\ Privacy Enhancing Technologies}, 2017(4).

\bibitem{barua2022blethreats}
A.~Barua, M.~A.~Al Alamin, M.~S.~Hossain, and E.~Hossain,
``Security and Privacy Threats for Bluetooth Low Energy in IoT and
Wearable Devices: A Comprehensive Survey,''
\emph{IEEE Open Journal of the Communications Society}, vol.~3, 2022.

\end{thebibliography}

\end{document}
